\chapter{System Model}\label{ch:system_model}

The model is organized around a deliberate resource separation. Communication
is a network-wide cooperative service, whereas sensing is accounted for through
a distinct, known covariance for each target. The AP--target association will
therefore answer a narrow physical question---which APs may allocate dedicated
sensing power to a target---without artificially restricting the communication
beamforming support. The following sections introduce this separation in the
same order in which it enters the optimization problem: network resources,
robust communication protection, and target-specific sensing accuracy.

\section{Network and Signalling Model}\label{sec:network_model}

We consider a downlink cell-free integrated sensing and communication (ISAC)
network in which a central processing unit coordinates $M$ distributed access
points (APs).  Each AP has $N_t$ transmit antennas and jointly serves $K$
single-antenna communication users while illuminating $P$ sensing targets.  The
AP, user, and target index sets are denoted by $\mathcal{M}=\{1,\ldots,M\}$,
$\mathcal{K}=\{1,\ldots,K\}$, and $\mathcal{P}=\{1,\ldots,P\}$,
respectively.  A subset $\mathcal{P}^{\mathrm{active}}\subseteq\mathcal{P}$
contains the targets that require sensing service in the current scheduling
interval.  The model is instantaneous within one coherence block; target-state
prediction across blocks is outside the scope of the resource-allocation
problem.

\begin{figure}[htbp]
    \centering
    \includegraphics[width=\linewidth]{../../paper_results_and_figures_v1_0/figures/archi.png}
    \caption{Cell-free ISAC architecture. The CPU coordinates distributed APs
    through fronthaul/control links (magenta). APs jointly serve users
    through globally cooperative communication beamforming (blue), whereas
    only associations with $b_{mp}=1$ authorize the target-specific dedicated
    sensing covariance $\mats{S}_p$ (orange). The generic-index illustration
    illustrates an inactive association, $b_{3,2}=0$, which denies AP 3
    dedicated sensing power for target 2 only; it does not prevent AP 3 from
    contributing to communication.}
    \label{fig:system_model}
\end{figure}

Let $\vect{w}_{m,k}\in\mathbb{C}^{N_t}$ be the communication beamformer from
AP $m$ to user $k$.  To make sensing participation explicit, a distinct
zero-mean sensing waveform is assigned to every target.  Its network-wide
stacked covariance is $\mats{S}_p\succeq\mats{0}$, where the stacked transmit
dimension is $N=MN_t$.  Define the stacked communication beamformer
\begin{equation}
    \vect{w}_{k}\triangleq[\vect{w}_{1,k}^{\T},\ldots,
    \vect{w}_{M,k}^{\T}]^{\T}\in\mathbb{C}^{N},
\end{equation}
and its covariance $\mats{W}_k=\vect{w}_k\vect{w}_k^{\H}$.  The total
transmit covariance is therefore
\begin{equation}\label{eq:total_covariance}
    \mats{R}_X=\sum_{k\in\mathcal{K}}\mats{W}_k+
    \sum_{p\in\mathcal{P}}\mats{S}_p.
\end{equation}
Let $\mats{E}_m\in\mathbb{R}^{N\times N}$ select the $N_t$-antenna block of
AP $m$.  The per-AP hardware constraint is
\begin{equation}\label{eq:ap_power_model}
    \Tr(\mats{E}_m\mats{R}_X)\leq P_{\max},\qquad \forall m\in\mathcal{M}.
\end{equation}

\section{Robust Communication Model}\label{sec:robust_comm_model}

The nominal stacked downlink channel of user $k$ is
$\hat{\vect{h}}_k\in\mathbb{C}^{N}$.  Imperfect channel state information
(CSI) is represented by the norm-bounded uncertainty model
\begin{equation}\label{eq:csi_uncertainty}
    \vect{h}_k=\hat{\vect{h}}_k+\Delta\vect{h}_k,\qquad
    \|\Delta\vect{h}_k\|_2\leq\epsilon_h\|\hat{\vect{h}}_k\|_2.
\end{equation}
Let $\mats{S}_{\Sigma}\triangleq\sum_{p\in\mathcal{P}}\mats{S}_p$ denote the
aggregate dedicated sensing covariance. It is treated as interference at a
communication user. This conservative and implementable assumption avoids
requiring receiver-side cancellation of the sensing waveforms. Thus the
worst-case downlink SINR requirement is
\begin{equation}\label{eq:robust_sinr_system}
 \min_{\|\Delta\vect{h}_k\|_2\leq\epsilon_h\|\hat{\vect{h}}_k\|_2}
 \frac{|(\hat{\vect{h}}_k+\Delta\vect{h}_k)^{\H}\vect{w}_k|^2}
 {\displaystyle\sum_{j\ne k}|(\hat{\vect{h}}_k+\Delta\vect{h}_k)^{\H}\vect{w}_j|^2
 +(\hat{\vect{h}}_k+\Delta\vect{h}_k)^{\H}
 \mats{S}_{\Sigma}(\hat{\vect{h}}_k+\Delta\vect{h}_k)+\sigma_c^2}
 \geq\gamma_k,\qquad \forall k\in\mathcal{K},
\end{equation}
where $\sigma_c^2$ is the communication-noise power and $\gamma_k$ is the
required SINR of user $k$.

\section{Target-Specific Sensing Model}\label{sec:sensing_model}

The binary variable $b_{mp}\in\{0,1\}$ indicates whether AP $m$ is admitted
to the sensing cluster of target $p$.  It authorizes dedicated sensing power
only; it does \emph{not} restrict the globally cooperative communication
beamformers.  For each active target, exactly $N_{\mathrm{req}}$ APs are
admitted:
\begin{equation}\label{eq:cluster_cardinality_model}
    \sum_{m\in\mathcal{M}}b_{mp}=N_{\mathrm{req}},\qquad
    \forall p\in\mathcal{P}^{\mathrm{active}}.
\end{equation}
The target-specific association gate is
\begin{equation}\label{eq:dedicated_sensing_gate}
    \Tr(\mats{E}_m\mats{S}_p)\leq P_{\max}b_{mp},
    \qquad\forall m\in\mathcal{M},\ p\in\mathcal{P}.
\end{equation}
Consequently, an unselected AP cannot allocate sensing power to target $p$,
whereas a selected AP is not forced to consume a minimum amount of power.  In
particular, \eqref{eq:dedicated_sensing_gate} is not a gate on communication
power and must be read jointly with the total hardware limit
\eqref{eq:ap_power_model}.

Let $\vect{g}_p\in\mathbb{C}^{N}$ be the effective sensing steering channel
of target $p$.  The target-specific sensing SINR is
\begin{equation}\label{eq:sensing_sinr_model}
    \mathrm{SINR}_{S,p}=
    \frac{\vect{g}_p^{\H}\mats{S}_p\vect{g}_p}{\sigma_s^2}
    \geq\gamma_S^{\mathrm{PoD}},\qquad \forall p\in\mathcal{P},
\end{equation}
where $\sigma_s^2$ is the sensing-noise power and
$\gamma_S^{\mathrm{PoD}}$ is the detection-related threshold.

For a target state $\boldsymbol{\theta}_p\in\mathbb{R}^{N_\theta}$, let
$\mats{D}_p=\partial\vect{g}_p/\partial\boldsymbol{\theta}_p
\in\mathbb{C}^{N\times N_\theta}$ be the effective response derivative.
Under the equivalent MISO sensing observation model, the data Fisher
information matrix (FIM) is
\begin{equation}\label{eq:fim_system}
    \mats{J}_p^{\mathrm{data}}(\mats{S}_p)=
    \frac{2}{\sigma_s^2}\Re\!\left\{
    \mats{D}_p^{\H}\mats{S}_p\mats{D}_p\right\}.
\end{equation}
Only the known waveform covariance $\mats{S}_p$ receives sensing-information
credit.  Communication covariances and waveforms dedicated to other targets
are not counted in \eqref{eq:fim_system}.  Tracking accuracy is constrained
by the matrix PCRB/CRB requirement
\begin{equation}\label{eq:pcrb_system}
    \Tr\!\left((\mats{J}_p^{\mathrm{data}})^{-1}\right)
    \leq\Gamma_{\mathrm{Track},p},\qquad\forall p\in\mathcal{P}.
\end{equation}
For $N_\theta=1$, this reduces to the scalar lower bound
$J_p^{\mathrm{data}}\geq1/\Gamma_{\mathrm{Track},p}$; for multi-dimensional
states, the full matrix form is retained.  The equivalent semidefinite form
used to solve \eqref{eq:pcrb_system} is developed in Chapter~\ref{ch:solution_method}.

\paragraph{Design implication.}
The model couples the three services through the common AP-power budget but
keeps their roles distinguishable: communication is protected for every
admissible CSI error, while a target receives sensing credit only from its own
authorized covariance. Chapter~\ref{ch:solution_method} turns these
requirements into the power-minimization problem and identifies the robust,
matrix, rank, and binary structures that prevent a direct convex solution.
